\documentclass[10pt]{article}
\usepackage[margin=20mm]{geometry}
\usepackage{amsmath,amssymb,mathtools,fontspec,unicode-math,xurl,hyperref}
\IfFontExistsTF{DejaVu Serif}{\setmainfont{DejaVu Serif}}{\setmainfont{Latin Modern Roman}}
\setmathfont{latinmodern-math.otf}
\setlength{\parindent}{0pt}\setlength{\parskip}{6pt}\sloppy
\title{Klein–Gordon Decay on Reissner–Nordström:
Review of the Acquired Preprint}
\author{AI-assisted technical review}
\date{}
\begin{document}
\maketitle
Abstract: This review concerns the exact acquired arXiv version of the paper by Shlapentokh-Rothman and Van de Moortel. Four core dependencies fail as written: the reference spectral operator is assigned a domain that excludes every nonzero eigenfunction; the threshold-frequency freezing is incompatible with the claimed uniformity at the event horizon; and an Airy-to-oscillatory connection coefficient has the opposite leading phase from the one required by its own Wronskians; and a quadratic source quantity is used to bound a linear Fourier response. Each finding has a bounded analytic contradiction and is classified Major. The corrected Dirichlet operator, a modulated horizon amplitude, a consistent change of matching phase, and a source norm of the correct homogeneity offer distinct repair routes. The last also requires renewed checks of frequency weights and the downstream integral estimates. These calculations do not disprove the final physical decay theorem after such repairs. Several further normalization and scope issues are recorded separately. Scientific adjudication remains pending; current focused independent acceptance is documented in the companion provenance, without human-review or publication claims.\par
\section*{Source identity and theorem scope}
The reviewed work is Polynomial time decay for solutions of the Klein–Gordon equation on a subextremal Reissner–Nordström black hole, by Yakov Shlapentokh-Rothman and Maxime Van de Moortel [paper], associated with DOI https://doi.org/10.1215/00127094-2025-0020. The actual scientific source is the 76-page arXiv:2401.00048v1, submitted on 29 December 2023 and internally dated 2 January 2024. Its matching author TeX was read as source, not executed. The publisher PDF was not obtained through the attempted authorized route. The journal metadata lists a substantially longer, 134-page article; no equivalence with that version is asserted. The original assigned batch membership is retained. All locations below use one-based PDF pages, matching printed pagination on the inspected pages.\par
The background has fixed \(M>0\), \(0\leq |e|<M\), and \(m>0\). Theorem 1.2, page 2, claims pointwise decay on \(r_+\leq r\leq R\), including the horizon, for smooth compactly supported initial data. Its unconditional loss is \(1/23\); an arbitrarily small positive loss is conditional on the stated exponent-pair conjecture. The fixed-angular-mode asymptotic in Theorem 1.1 is imported from the authors’ previous work and is not a new uniform-in-mode estimate. The quasimodes of Theorem 1.4 move to \(r\gtrsim L^2\). Their logarithmic obstruction for polynomially localized data in moving large-radius regions does not contradict a decay statement for fixed compact observation regions and sufficiently localized data.\par
The proof chain runs from the separated radial equation through high- and intermediate-frequency estimates, three turning-point connections, a resonant Green formula, stationary phase, exponential sums, and angular Sobolev summation. The findings below concern specific links in that chain. Native AI helpers performed bounded section and source-page checks, whose roles, exact coverage and limitations are retained in the accompanying records. They are not human reviewers or independent acceptance agents.\par
\section*{The reference spectral operator}
Problem 1 (Two traces vanish at the regular endpoint).  Definition 4.2 and Lemma 4.7, page 20, assign the half-line Schrödinger operator \(Q\) the domain \(H^2_0((x_0,\infty))\) and assert self-adjointness. The preceding sentence explicitly defines \(H^k_0\) as the \(H^k\)-closure of compactly supported smooth functions. With this definition the domain is too small.\par
\subsection*{An admissible contradiction}
Choose \[M=m=1,\qquad e=0,\qquad x_0=2,\qquad
 \mathcal E_1=\mathcal E_2=0,\] and any sufficiently large integer \(L\). All the printed potential bounds are satisfied. The actual operator is then \[Qf=-f''+\left(-\frac2x+\frac{L(L+1)}{x^2}\right)f,\qquad
 D(Q)=H^2_0((2,\infty)).\] Continuity of the two endpoint trace maps on \(H^2\) implies \[f(2)=f'(2)=0\qquad(f\in D(Q)).\] If \(Qf=\lambda f\), the regular second-order equation \[f''=\left(-\frac2x+\frac{L(L+1)}{x^2}-\lambda\right)f\] has zero Cauchy data at \(2\). Uniqueness gives \(f\equiv0\). Thus this operator has no eigenfunctions, in direct contradiction to Lemma 4.8, page 21, which asserts infinitely many negative eigenvalues tending to zero.\par
Self-adjointness can also be refuted without any spectral calculation. Let \(g(x)=e^{-(x-2)}\). Both \(g\) and \[-g''+\left(-\frac2x+\frac{L(L+1)}{x^2}\right)g\] belong to \(L^2((2,\infty))\). Integration by parts against \(f\in H^2_0\) has no endpoint term because both traces of \(f\) vanish. Hence \(g\in D(Q^*)\). But \(g(2)=1\), so \(g\notin D(Q)\). Therefore \(Q\ne Q^*\).\par
\subsection*{Core impact and the repair boundary}
This is a Major defect in a core spectral dependency. The stated operator cannot provide the eigenfunctions, min–max principle, eigenvalue sequence or spectral resolvent used in Section 4.2 and the later frequency decomposition. In particular, the quantitative eigenvalue comparison of Lemma 4.11, pages 22–23, is not an estimate for the operator actually defined.\par
The source also gives strong evidence for the intended correction: the shifted sine-transform argument in Lemma 4.8 describes the Dirichlet realization. Its domain is \[D(Q_D)=H^2((2,\infty))\cap H^1_0((2,\infty)),\] which imposes \(f(2)=0\) but not \(f'(2)=0\). Equivalently, one can specify the Friedrichs operator associated with the \(H^1_0\) quadratic form. For this bounded real potential the Dirichlet realization is self-adjoint. It can have the Coulomb-type negative spectrum. The present counterexample does not refute the spectral estimates for that corrected operator; they must be stated and used for it consistently.\par
A related definition must be repaired in the frequency subdivision. Where distance to the spectrum is intended, the quantity is \[\inf_n|m^2-\omega^2+\lambda_n|.\] A signed infimum instead equals \(m^2-\omega^2+\lambda_1\) for an increasing eigenvalue sequence. Its sublevel set contains an interval of algebraic width, rather than only the intended small neighborhoods of the individual eigenvalues. The later lemmas use the absolute distance, whereas Section 6 uses the small-measure decomposition. This is recorded as an additional consistency requirement within the same spectral dependency, not as a separate constructed resolvent-growth counterexample.\par
\section*{Uniformity of the threshold Green formula at the horizon}
Problem 2 (Freezing an oscillatory horizon factor is not uniform).  Theorem 5.1(4), page 32, equation (5.10), evaluates the regular Green coefficients at \(\omega=m\). They contain \(u_H(m,s)\), while the error bounds explicitly hold for every \(s\leq s_{\max}\). The horizon normalization (5.5), page 31, makes these two requirements incompatible.\par
\subsection*{An exact source within the theorem’s hypotheses}
Fix a subextremal background and one sufficiently large \(L\); keep \(L\) fixed in the following threshold limit. Choose a real smooth cutoff \(\chi\) equal to \(1\) for \(s\leq0\) and \(0\) for \(s\geq1\). Define \[u(\omega,s)=\chi(s)u_H(\omega,s),\qquad
 H(\omega,s)=\chi''(s)u_H(\omega,s)
                 +2\chi'(s)\partial_su_H(\omega,s).\] Because \(u_H''=Vu_H\), one has \(u''-Vu=H\) exactly. Both \(H\) and \(\partial_\omega H\) are supported in \([0,1]\), and hence satisfy the theorem’s source support condition for any \(s_H>1\). On this fixed interval their norms are finite uniformly for \(m/2\leq\omega\leq m\), by the horizon Volterra construction away from zero frequency. In particular the printed quantity \(N[H]\) is finite at this fixed \(L\).\par
The solution is outgoing at the horizon and identically zero near infinity, an admissible decaying end condition for the inhomogeneous linear problem. A unit nonzero outgoing amplitude cannot be imposed on every inhomogeneous solution: it would exclude zero forcing and contradict the linear Green formula itself. If desired, a smooth frequency cutoff and conjugate reflection at negative frequencies give a real smooth physical source. The local theorem does not require time-compact support.\par
For \(s<0\), the two \(u^{\rm reg}_2\) integrals in (5.10) vanish. The two remaining regular terms depend on \(s\) only through \(u_H(m,s)\). Consequently the theorem implies \[u_H(\omega,s)=B_L(\omega)u_H(m,s)+R_L(\omega,s),
 \qquad
 \sup_{s<0}|R_L(\omega,s)|\leq C_LN[H]k^{-1/2},\] for some complex scalar \(B_L(\omega)\). Here \[k^{-2}=\frac{m^2-\omega^2}{(Mm^2)^2}.\] The constant \(C_L\) may be very large, but is finite and independent of \(\omega\) at fixed \(L\). Indeed the theorem states \[d_L:=1-|\Gamma(m,L)|>0,
 \qquad |1+\Gamma(m,L)e^{2i\vartheta}|\geq d_L.\] Thus the denominator and all bounded trigonometric factors can be absorbed into \(C_L\). No lower bound uniform in \(L\) is needed.\par
\subsection*{The phase-circle contradiction}
Equation (5.5) gives, for each fixed \(\omega\), \[u_H(\omega,s)=e^{-i\omega s}(1+o(1)),
 \qquad u_H(m,s)=e^{-ims}(1+o(1))
 \quad(s\longrightarrow-\infty).\] For every \(\omega\ne m\) and every \(B\in\mathbb C\), \[\limsup_{s\to-\infty}|u_H(\omega,s)-Bu_H(m,s)|
 =\max_{|z|=1}|z-B|=1+|B|\geq1.\] The phase \(e^{i(m-\omega)s}\) traverses the unit circle infinitely often. The two \(o(1)\) terms do not change this limit superior. But the claimed upper bound tends to zero as \(\omega\uparrow m\), since \(k\to\infty\). This is a contradiction with the theorem’s actual quantifiers.\par
\subsection*{Propagation, contrary evidence and correction}
The proof in Section 5.12, page 56, freezes all regular coefficients and invokes a frequency derivative bound. The unmodulated factor has a derivative whose leading term is \(-is e^{-i\omega s}\), which is not uniform on the stated half-line. Section 6 retains frozen factors multiplied by \(e^{i\omega p(s)}\). Near the horizon the residual factor is \(e^{i(\omega-m)s}\); multiplication by the horizon-penetrating Fourier phase does not remove the error.\par
This is a Major failure of the central uniform resolvent representation used for the announced decay up to the horizon. It is unaffected by Wronskian signs, by a correction of the action’s \(\pi\) normalization, or by keeping the source away from the horizon; the constructed source already has that support.\par
The appropriate regular object is \(e^{i\omega s}u_H(\omega,s)\), which tends to \(1\) at the horizon. The source uses such modulated derivatives elsewhere, providing evidence for an intended repair. One must retain \(e^{-i\omega s}\) and freeze only the regular amplitude, then re-establish both Green terms and the inverse-Fourier remainder estimates. Restricting \(s\) to a fixed bounded interval is a narrower repair that does not prove the claimed horizon-uniform result. No counterexample to the final physical decay rate is asserted.\par
\section*{The second turning-point phase}
Problem 3 (The inverse connection has the wrong leading phase).  Corollary 5.13, page 48, equation (5.45), gives a diagonal connection coefficient with phase \(e^{-i\pi/4}\). The forward expansion of Corollary 5.8, page 44, and the exact Wronskian normalization (5.36), page 42, require \(e^{+i\pi/4}\).\par
\subsection*{A bounded Wronskian calculation}
Write \[A=u_{Ai_2},\quad B=u_{Bi_2},\quad
 c=(Mm^2)^{1/2}\pi^{-1/2}\mathbb L^{-2/3},
 \quad F_+=A-iB,\quad F_-=A+iB.\] The paper’s convention is \(W(f,g)=f'g-fg'\). Its exact normalization is \(W(A,B)=c^2\), hence \[W(F_+,F_-)=2ic^2.\] The oscillatory solutions satisfy \(W(B_+,B_-)=2i\). Corollary 5.8 and the proof of Corollary 5.13 give \[F_-=c\,e^{i\pi/4}B_-+cR_-,
 \qquad
 W(B_+,R_-)=O(e^{-\kappa_+A_0}+\mathbb L^{-1}),\] where \(A_0\) denotes the paper’s large matching parameter, distinguished here from the Airy solution \(A\).\par
In the exact expansion \(B_+=b_{++}F_++b_{+-}F_-\), taking a Wronskian with \(F_-\) gives \[b_{++}
  =\frac{W(B_+,F_-)}{W(F_+,F_-)}
  =c^{-1}\left[e^{i\pi/4}
                  +O(e^{-\kappa_+A_0}+\mathbb L^{-1})\right].\] The displayed (5.45) instead asserts \[cb_{++}=e^{-i\pi/4}+O(e^{-\kappa_+A_0}).\] The normalized leading discrepancy is exactly \[|e^{i\pi/4}-e^{-i\pi/4}|=\sqrt2.\] Take, for example, the admissible fixed background \(M=m=1,e=0\). The source’s convention allows the matching parameter \(A_0\) large and then \(\mathbb L\) sufficiently large relative to it, with \(k\) in the stated threshold regime. The unqualified error constants in (5.45) depend on the background, not on \(A_0\); neighboring estimates explicitly use an \(A_0\) subscript when that dependence is allowed. Thus the errors can be made arbitrarily small, while the discrepancy remains \(\sqrt2\). This calculation isolates the actual coefficient and does not depend on a choice of observation point where two oscillations might cancel.\par
\subsection*{Why the correction reaches the resonant formula}
The same phase is carried into Lemma 5.16, page 50. With the source’s real-basis definitions, \[w_{\rm even}=\Re(e^{i\pi/4}w_{2,+}),\qquad
 w_{\rm odd}=-\Im(e^{i\pi/4}w_{2,+}),\] the corrected leading connection gives \[w_{\rm even}\sim c^{-1}B,\qquad
 w_{\rm odd}\sim-c^{-1}A,\] rather than the printed \(c^{-1}A,c^{-1}B\). The leading sine and cosine roles in Proposition 5.21, page 54, therefore rotate. The subsequent Wronskian and narrow-resonance formula must be recomputed with one consistent phase convention. A fixed \(\pi/2\) shift cannot be absorbed into the stated \(\phi_3=O(k^{-1})\).\par
This is a Major error in the central matching calculation. The bounded repair begins by conjugating the leading diagonal phases in (5.45), followed by propagation through the real basis, connection coefficients and resonant phase. A global basis convention could also be changed, but would require changing the already fixed Airy normalization and all dependent formulas consistently. The conserved flux calculation supplies useful contrary evidence: with a consistent basis, it still enforces the correct modulus gap \(1-|\Gamma|>0\). The existence of a gap and the general tunneling mechanism are not refuted by the phase error. Nor does this calculation alone determine that the final polynomial decay exponent changes.\par
\section*{Homogeneity of the source-to-solution estimate}
Problem 4 (A quadratic source quantity cannot bound a linear response).  Proposition 6.6, page 62, defines its quantity \(\mathcal D[H]\) as a supremum of weighted squared source integrals, without a square root. Its Fourier-response bound is linear in \(u\) but multiplies \(\mathcal D[H]\) once. There is no unit-source normalization or lower bound on \(\mathcal D\). The same mismatch occurs in Proposition 3.2, page 12, and Lemma 6.7, page 62.\par
\subsection*{A nonzero response satisfying the source conditions}
Fix the background, the spatial cutoff \(S_0\), the required large weight exponent, \(R>r_+\), and \(\check\delta>0\). Choose \(s_0<S_0\) with \(r(s_0)\in(r_+,R)\), and \[\chi\in C_c^\infty((-\infty,S_0)),\quad \chi(s_0)=1,\qquad
 0\ne\eta\in C_c^\infty((m+1,m+2)),\quad\eta\geq0.\] Take a fixed finite \(t_0\geq1\) such that \(T_0=t_0+p(s_0)\geq1\). Choose a sufficiently large fixed \(L_0\) so that \(t_0,T_0<\exp(\check\delta\bar L_0\log\bar L_0)\). Set \[\begin{split}
 u_0(\omega,L,s)&=1_{\{L=L_0\}}
                 e^{i\omega t_0}\eta(\omega)\chi(s),\\
 H_0(\omega,L,s)&=(\partial_s^2-V)u_0(\omega,L,s).
 \end{split}\] The equation holds exactly. Spatial support is compact and below \(S_0\); frequency support is compact and away from \(0,m\). All required derivatives are smooth. Only one finite angular mode is used, and \((r-r_+)^{-1}\) is bounded on the source support. Thus \(0<\mathcal D[H_0]<\infty\): if \(H_0\) were zero, uniqueness for the homogeneous ODE and the vanishing solution outside its compact support would force \(u_0=0\).\par
Both end amplitudes vanish, so \(u_0\) lies in the linear outgoing subspaces used by the paper’s Green formula. This convention is material: if the earlier symbol \(\sim\) were instead read as requiring unit amplitude at both ends for every inhomogeneous solution, it would conflict with that linear formula, the zero source and the final application to arbitrary initial data. The present construction uses the source-to-solution interpretation actually required in those applications, not a hidden unit-amplitude restriction.\par
At the selected observation point, \[\int_0^\infty e^{-iT_0\omega}e^{i\omega p(s_0)}
       u_0(\omega,L_0,s_0)\,d\omega
   =\int_0^\infty\eta(\omega)\,d\omega=:c_0>0.\] For \(u_\varepsilon=\varepsilon u_0\) and \(H_\varepsilon=\varepsilon H_0\), the source support and every fixed parameter remain unchanged, but \[\mathcal D[H_\varepsilon]=\varepsilon^2\mathcal D[H_0].\] The indicator term in Proposition 6.6 is zero at \(t_0\). Its remaining bound would therefore imply \(\varepsilon c_0\leq C\varepsilon^2\mathcal D[H_0]\), with \(C\) finite and independent of \(\varepsilon\). Letting \(\varepsilon\downarrow0\) gives a contradiction.\par
\subsection*{Actual dependence and limits of the repair}
The final proof on pages 68–69 explicitly invokes Proposition 6.6 and Lemma 6.15 for the time-cutoff source \(F=2\nabla\chi\cdot\nabla\phi+(\Box\chi)\phi\). For actual compactly supported smooth initial data, a fixed smooth cutoff and finite propagation preserve linear scaling: \(\phi,F,H,u\) scale by \(\varepsilon\), whereas \(\mathcal D\) scales by \(\varepsilon^2\). A nonzero compact exterior single-mode solution at a fixed finite time can be chosen by finite-time Cauchy evolution. Its cutoff source is supported in a compact radial region during the transition, so its Fourier source has finite \(\mathcal D\). The outgoing and Fourier-inversion interpretation here is the one imported in the paper’s final proof; that external scattering theorem is not independently rederived in this review. The explicit construction above already refutes the stated general-source estimate. The final claim that the same quadratic \(\mathcal D\) is bounded by a first-power Sobolev norm also fails under large-amplitude scaling whenever \(\mathcal D>0\).\par
This is a Major defect in a core estimate, grouped as one problem across its repeated uses. Proposition 3.1 has a squared left side and supplies contrary evidence for the intended homogeneity. Taking a square root of the quadratic quantity restores necessary linear scaling, but does not certify the complete repair: frequency and angular weights and all subsequent estimates must be checked again. For example, the displayed \(\omega^8\) bound in Proposition 3.1, combined only with the \((1+\omega^{10})\) weight before taking square roots, gives a coarse \(\omega^{-1}\) derivative bound, insufficient by itself for the high-frequency \(L^1_\omega\) step in Lemma 6.7. Stronger weights or finer estimates are needed there. The final theorem has a first-power data norm and is not itself refuted by this scaling argument.\par
\section*{Checks that survive and further bounded corrections}
\subsection*{The exponential-sum mechanism}
The compensation between a narrow resonance and its prefactor survives an independent algebraic check. Let \(a\asymp1-r\), \(\rho=-\Gamma\), \(r=|\rho|<1\), and \(|1-\rho|\lesssim a\). If \[|b_q|\lesssim q^{1/3},\quad
 |b_{q+1}-b_q|\lesssim q^{-2/3},\quad
 |S_q|\leq Kq^{2/3},\] Abel summation gives \[a\left|\sum_{q\leq N}\rho^q b_qz_q\right|
 \lesssim K\left[
 aNr^N+a\sum_{q\geq1}r^q
       +a^2\sum_{q\geq1}qr^q\right]\lesssim K.\] Thus the large angular parameter need not destroy the sum estimate once the correct phase and amplitude hypotheses have been proved.\par
The dyadic exponent arithmetic also has a positive margin. For the two exponent pairs used in Appendix C, the competing exponents are \[E_1=\frac{1+884\alpha}{2745},
 \qquad E_2=\frac{13}{252}-\frac{4\alpha}{63}.\] Optimizing gives \(\alpha=3937/29632\) and the common value \(959/22224\). The difference from \(1/23\) is \(167/511152>0\), leaving room for sufficiently small losses. This checks the stated arithmetic mechanism; it neither proves the exponent-pair conjecture nor repairs the preceding analytic dependencies.\par
\subsection*{Definitions and local normalizations}
The coordinate displays on pages 10–11 need correction. The metric with \(f=1-2M/r+e^2/r^2\) requires \(ds/dr=f^{-1}\). For \(M=1,e=0\), the explicit printed coordinate is \(s=r+2\log[r(r-2)]\), so \[s'=\frac1f+\frac2r,
 \qquad s'(4)=\frac52\ne2=\frac1{f(4)}.\] The displayed radial operator \(f\partial_r(f^{-1}\partial_r)\) also differs from \(f\partial_r(f\partial_r)\). Applied to \(\psi=r\) at \(r=4\), the values are respectively \(-1/4\) and \(1/16\). The later analysis largely uses the intended tortoise-coordinate comparabilities. An additive normalization of the correct coordinate can preserve the displayed leading horizon constant \(C_+\). These observations establish the local errors, without separately refuting the subsequent decay inequalities.\par
The interpolation conditions on \(P\) in Definition 3.1, page 13, also need a low-frequency qualification: when \(M=1,e=0,m=1/2,\omega=1\), the asserted \(|\partial_s\log P|\lesssim S_\xi^{-1}\omega\log\omega\) forces a constant transition. Its prescribed outer value is \(3/4+1/(2r(2S_\xi))<1\) for large \(S_\xi\), whereas its inner value is \(1\). A different bounded-frequency interpolation is needed. This does not by itself refute the later large-frequency estimate.\par
Appendix A, page 69, uses an \(I_{i\nu}-I_{-i\nu}\) difference with the opposite sign from the standard \(K_{i\nu}\) connection formula [dlmf]. Its small-argument sign agrees with that difference, but its positive large-argument leading term agrees with the standard \(K\). At \(\nu=x=1\), the difference displayed in the paper is approximately \(-0.2894280370\), while \(K_i(1)\) is approximately \(+0.2894280370\). The exact connection identity, rather than this numerical illustration, proves the sign conflict. A consistent basis-sign correction is required; the modulus flux identity is not independently refuted.\par
Other recorded local corrections include the extra \(\pi^{-1}\) in (5.9) relative to the action definition in Proposition 5.22; and the derivative coefficients of the leading \(A^{2/3}\) phase in Section 6. The corrected derivative profile still has the shape needed for the exponent-pair argument. These remaining points retain their bounded uses. The source-norm homogeneity defect has instead been isolated as the core Problem 4.\par
\subsection*{Quantifiers requiring clarification}
Theorem 1.2 suppresses dependence on the support radius of the initial data, although the proved resolvent bounds explicitly depend on the support cutoff \(s_H\) or \(S_0\). The final proof does not supply a bound uniform as that cutoff tends to infinity. The source should either fix the support radius and display this dependence, introduce suitable spatial weights, or supply a uniform argument. No complete black-hole counterexample to that uniformity has been established here; this remains a request for evidence.\par
Similarly, the final statement about every coordinate derivative needs the appropriate dependence of the data regularity order on the derivative order. The argument only calls \(N\) sufficiently large. It should not be read as supplying one fixed finite Sobolev order for unbounded derivative orders without further proof. The dependence of the frequency cutoff exponent on the requested small loss must also be aligned with the final Sobolev-order quantifier.\par
A further bounded question concerns the geometric-tail derivative estimate (6.23), page 67. This concerns the genuine tail of the geometric expansion (6.19), with the summation index \(q\) restored in the subsequent exponential factors. Those factors, including the displayed definition of \(F\) preceding (6.23), omit \(q\) as printed and do not literally describe the same tail. After correcting the source homogeneity in Problem 4, consider the model tail \[F(\theta)=A_2\frac{(re^{i\theta})^{N+1}}{1-re^{i\theta}},\] its exact derivative at \(\theta=0\), with \(A_2,r\) fixed, has size \[|A_2|r^{N+1}|\theta'|
 \left[\frac{N+1}{1-r}+\frac{r}{(1-r)^2}\right].\] When \(|A_2|\asymp a\), \(1-r\asymp a\), and \(Na\ll1\), this can have size \(|\theta'|/a\). The stated geometric bounds alone therefore do not give the printed pointwise derivative gain. An actual lower bound for the data-dependent coefficient at a resonant frequency was not established here, so this is an inference diagnostic, not a counterexample to the physical solution. Phase averaging may instead prove the needed integrated estimate, since \(a\int_0^{2\pi}|1-re^{i\theta}|^{-2}\,d\theta=O(1)\). That repaired argument remains to be supplied.\par
\section*{Assessment and verification limits}
The four Major findings concern the operator domain, the horizon-uniform resolvent error, the turning-point connection phase, and the nonhomogeneous source-to-solution estimate. Each is supported by an explicit admissible calculation, the exact source location, a check of contrary evidence, and a bounded repair statement. None requires the entire paper to be false.\par
The acquisition record includes the available author source archive, metadata and unsuccessful publisher route. No paper-specific computational supplement or code was identified in the available searches; this is an access and reproducibility limitation, not a mathematical defect. No author code or author TeX was executed. External Olver, spectral-theoretic, prior fixed-mode scattering and number-theoretic inputs were not all independently reproved. This focused acceptance does not certify a corrected complete argument.\par
Source, TeX, dependencies, logs, resulting PDF and rendered page images are bound by actual hashes in the accompanying records. Actual page-inspection receipts are separate from render-only records. Native assistance is attributed by task and coverage; it is not represented as human review or as an independently observed backend model. Scientific decision remains pending and the human-review array is empty.\par
Exact reviewed source SHA-256: `20a6097cff0ac353afa8b7cb63ac5c8d9b4bc152c507eb721bed52f02483a04b`.\par
9 Y. Shlapentokh-Rothman and M. Van de Moortel, Polynomial time decay for solutions of the Klein–Gordon equation on a subextremal Reissner–Nordström black hole, arXiv:2401.00048v1, submitted 29 December 2023, internal date 2 January 2024, exact acquired 76-page PDF. https://arxiv.org/abs/2401.00048v1. NIST Digital Library of Mathematical Functions, §10.27, equation 10.27.4, modified Bessel connection formula, consulted during this review. https://dlmf.nist.gov/10.27.E4.\par
\end{document}
